\documentclass[11pt]{article}
\usepackage[margin=1in]{geometry}
\usepackage[T1]{fontenc}
\usepackage{amsmath,amssymb,amsthm}
\usepackage[hidelinks]{hyperref}
\newtheorem{theorem}{Theorem}
\setlength{\emergencystretch}{3em}
\title{Finite polyannulus Picard vanishing, a nonalgebraizable surface, and bounded branch energy}
\author{Independent P16.7 research draft}
\date{8 October 2026; uncompiled and unverified}
\begin{document}
\maketitle
Unverified research draft, 8 October 2026. ONLY original P16.7 is active. Original completion: FALSE. The new finite base-field Picard bridge and bounded branch results are proved below under the explicitly stated published analytic foundations. No full intrinsic unbounded E1 equation theorem or admissible counterexample is claimed. Raw analytic mass is two. Saved TeX is uncompiled; ROOT retains compilation and acceptance.

\section*{Precise new partial results}
Retain $K=\mathbb C((s))$, $|s|=e^{-1/8}$, $t=s^8$.
\begin{theorem}[Finite Picard bridge]
Every finite product of closed discs, circles and nonzero-width closed annuli
with positive radii in $|K^\times|$ has zero analytic Picard group.
In particular $\operatorname{Pic}(P_m)=0$ for all integers $m\geq1$,
where $P_m=\{|s|^m\leq|z_j|\leq|s|^{-m},\ j=1,2\}$.
\end{theorem}
The complete proof in Section 1 applies Gruson only to pure disc/circle products
and uses coherent finite rational-cover gluing to fill each hole.
\begin{theorem}[Chosen base-field example]
The quotient $X=(\mathbb G_m^2)^{\rm an}/\langle(t,2),(3,t)\rangle$ is good,
strict, smooth, separated, proper and boundaryless. Its meromorphic field is $K$.
No separated finite-type $K$-algebraic space has analytification $X$.
\end{theorem}
The full compatible-frame, pole-divisor and algebraic-space proof is Section 2.
\begin{theorem}[Bounded branch pencils]
For the real analytic datum $\theta=d'd''(|x|^2/2)$ of raw mass two and the
specified ordinary locally smooth-positive bounded class $\mathcal C_\theta$,
the potentials $V_{a,N}=u_0+aD_N$, $0\leq a\leq1$, $N\geq1$ real, satisfy
\[
 E_\theta(V_{a,N})=-\frac16-\frac{a^3N}{3},\qquad
 \operatorname{MA}_\theta(V_{a,N})=(1-a^2)\operatorname{MA}_\theta(u_0)
      +a^2\operatorname{MA}_\theta(U_N).
\]
For $v_N=V_{N^{-1/3},N}-1/12$ one has
\[
 \sup_Xv_N=0,\quad E_\theta(v_N)=-\frac23,\quad
 v_N(1,1)=-\frac14-N^{2/3}.
\]
Consequently this normalized energy-bounded family fails uniform relative compactness.
If a positive raw-mass-two measure $\mu$ has finite bounded dual energy $M$,
put $C=M-[E_\theta(u_0)-\int u_0\,d\mu]$. Then for $N\geq1$ and $N>3C/2$,
\[
 \mu\{ |z_j-1|\leq e^{-N},\ j=1,2\}
 \leq (3C/2)^{2/3}N^{-2/3}.
\]
\end{theorem}
These statements do not establish compactness in a weaker topology,
an unbounded operator, a sufficient finite-energy criterion or the original equation theorem.
The complete bounded-current and ordinary-support proof is Section 3.

\section{n1. Filling annular holes closes the finite base-K Picard gate}

This is an unverified proof draft, not an execution certificate. Retain K=C((s)) with |s|=exp(-1/8). The result is Pic(P\_m,K)=0 for every integer m>=1, where P\_m,K is the product of the two closed annuli |s|\textasciicircum{}m<=|z\_j|<=|s|\textasciicircum{}(-m). In fact the argument proves the following slightly larger finite statement over this same K: every finite product of closed discs, circles, and closed annuli whose positive radii belong to |K*| has zero analytic Picard group. No arbitrary-extension quantifier is needed or claimed.

The source input is precisely Gruson, *Fibres vectoriels sur un polydisque ultrametrique*, Ann. Sci. ENS (4) 1 (1968), 45--89, IV.5, Theorem 1 (printed pp. 79--81). Its H is a completed tensor product of disc algebras T\_(1,r), r in Q tensor v(K*), and unit-circle algebras C. The conclusion includes SK\textasciicircum{}0(H)=0; SK\textasciicircum{}0 is the cokernel of Z->[K\textasciicircum{}0(H)] sending 1 to [H]. This theorem concerns discs and circles. It is never applied here with a wide-annulus factor. Its full immediate proof and the K\textasciicircum{}0 dependencies are recorded in the local source ledger. Only the K\textasciicircum{}0 conclusion is used, so its additional SK\textasciicircum{}1 calculation and its valued-excellence dependency are not needed for the inference below.

First, SK\textasciicircum{}0(H)=0 implies Pic(H)=0. The finite analytic line bundles on an affinoid correspond to finite projective modules of rank one. If P has rank one, its class in K\textasciicircum{}0(H) is n[H] for an integer n; evaluating rank at any prime gives n=1. Equality in the Grothendieck group of the direct-sum monoid of finite projectives means P direct-sum Q is isomorphic to H direct-sum Q for some finite projective Q. Taking top exterior powers gives P tensor det(Q) isomorphic to det(Q). Tensoring with the dual of det(Q) gives P isomorphic to H. Thus no cancellation assertion for arbitrary vector bundles is required. A circle of radius c in |K*| is carried to a unit circle by z->z/a with |a|=c. Discs are similarly rescaled to unit discs. This proves the zero-Picard assertion for products having no wide-annulus factors.

We use analytic G-sheaves for line bundles. Finite rational-domain covers support sheaf gluing; the affinoid correspondence and coherent descent are the Tate--Kiehl foundations, in the precise formulation of Van der Put, *Rigid analytic spaces* (1975--76), 2.14--2.16.3 and 3.4.5--6. The complete immediate proof of 2.16.3, including its explicit descent diagram, has been read; it cites Kiehl for the deeper affinoid descent theorem. These covers are not represented as ordinary open covers for a psh construction.

Here is the induction, which repairs the foundation without asserting arbitrary-base factoriality. Write D(R)=\{|z|<=R\}, D(r)=\{|z|<=r\}, A(r,R)=\{r<=|z|<=R\}, and C(r)=\{|z|=r\}, with 0<r<R in |K*|. Suppose all finite products with fewer than k wide-annulus factors, and any number of disc or circle factors, have zero Picard group. Let Z=A(r,R) times Y have k such factors. The product D(R) times Y has k-1; the overlap C(r) times Y has k-1 as well.

Take ANY analytic line bundle L on Z. Its restriction to C(r) times Y is trivial by induction. Choose one frame there. On D(r) times Y take the trivial bundle and identify its frame with that chosen frame on the overlap. The two analytic domains A(r,R) times Y and D(r) times Y cover D(R) times Y; they are rational domains, their intersection is exactly C(r) times Y, and they form a finite admissible affinoid cover. For example, the two inequalities |z/a|<=1 and |a/z|<=1, with |a|=r, give the disc and annular sides; the denominator z together with the nonzero constant a generates the unit ideal. Their products with Y remain rational domains. The chosen isomorphism and its inverse give descent data; the two-member cocycle condition is automatic. Ordinary sheaf gluing on this analytic G-cover produces an invertible sheaf L' on D(R) times Y. Local freeness and coherence descend by the stated affinoid theorem. By induction L' is trivial, so its restriction L is trivial. This proves the induction for every k.

For the exact two-variable case one may check the order without any induction shorthand. First prove Pic(C\_1 times A\_2)=0 by extending across the second hole: the overlap C\_1 times C\_2 and the filled space C\_1 times D\_2 are pure circle/disc products. Next extend an arbitrary L on A\_1 times A\_2 to D\_1 times A\_2 using that boundary result. Finally extend to D\_1 times D\_2 using the pure product D\_1 times C\_2 on the second overlap. The bidisc is covered by Gruson. In particular, the delicate intermediate D\_1 times A\_2 is proved by filling, not by applying his theorem to a wide-annulus base.

Taking r=|s|\textasciicircum{}m, R=|s|\textasciicircum{}(-m) establishes the exact base-field gate F\_K for every m. No topological contractibility argument, formal-model extension assertion, pro-homotopy limit, or regular-factorial-base restricted-variable theorem is used. This is a complete finite Picard bridge under the stated standard affinoid foundations, not the original finite-energy equation. It makes no assertion about arbitrary radii outside |K*| or arbitrary complete extensions L/K.

Artifacts: notes/n1.md and P16.7-standalone-paper.tex. Source units and physical-reading distinctions are in source-application-ledger.md and sources-LOCALONLY/. Own finite gluing-combinatorics checks are rerun with `python scripts/diagnostics.py`; they illustrate the induction and do not prove analytic descent. Support: none; no supplied note is a premise.


\section{n2. Base-field nonalgebraization via compatible frames and all pole divisors}

Retain K=C((s)), |s|=exp(-1/8), t=s\textasciicircum{}8. Let T=(G\_m\textasciicircum{}2)\textasciicircum{}an, lambda\_1=(t,2), lambda\_2=(3,t), Lambda=<lambda\_1,lambda\_2>, and X=T/Lambda. With the finite Picard bridge n1, the following argument no longer assumes F\_K. It excludes EVERY separated finite-type K-algebraic space algebraization of X, not just projective schemes or algebraic group uniformizations. This is an unverified mathematical draft using the explicitly named analytic foundations; it does not solve P16.7.

Write x=(-log|z\_1|,-log|z\_2|). Multiplication by lambda(n)=lambda\_1\textasciicircum{}(n\_1)lambda\_2\textasciicircum{}(n\_2) sends x to x+n. Inverse images of open rectangles of widths less than one have disjoint nonidentity deck translates. Their quotient maps are ordinary open analytic charts, glued by locally constant deck transformations satisfying the cocycle. This constructs X. The applied open-quotient gluing statement and its entire immediate proof are Conrad--Temkin, *Non-archimedean analytification of algebraic spaces*, author version dated 24 February 2009, Lemma 4.2.3. The action relation in T times\_K T is a locally finite union of closed translation graphs: on any pair of relatively compact subsets of T only finitely many n can occur, since n lies in a bounded difference of tropical-coordinate ranges. Thus its map to T times\_K T is a closed immersion. Closed diagonals descend through the local isomorphism atlas; Conrad--Temkin 4.2.4 and its adjacent closed-immersion descent proof justify this in the analytic category, beyond mere Hausdorffness. X is separated.

The compact strict polyannulus x\textasciicircum{}(-1)([0,1]\textasciicircum{}2) surjects onto X. Goodness and strictness can also be seen pointwise: choose a center in (1/8 Z)\textasciicircum{}2 within 1/16 in each coordinate of a lift; the closed rectangle of width 1/2 about it lies in the open rectangle of width 3/4, which is a quotient chart. Its inverse image is a strict affinoid neighborhood of the point. Local charts are ordinary opens of the smooth boundaryless torus. Consequently X is good, strict, smooth, of pure dimension two, boundaryless and compact; with separatedness these are precisely analytic properness. T is connected: its nested centered closed polyannuli have integral Laurent algebras and nonempty intersections. Hence X is connected. Only this base-K geometry is asserted here.

We first prove Pic(T)=0, including compatibility along the noncompact exhaustion. Put a=1/8 and P\_m=\{|s|\textasciicircum{}m<=|z\_i|<=|s|\textasciicircum{}(-m)\}. By n1 choose a frame e\_m of a given line bundle on P\_m. Every unit f on P\_m has an inner estimate as follows. Its convergent Laurent expansion is sum b\_beta z\textasciicircum{}beta. At the Gauss seminorm of log-radii y in [-ma,ma]\textasciicircum{}2, log|f| is the maximum of log|b\_beta|+beta dot y. It and the corresponding function for f\textasciicircum{}(-1) are convex and sum to zero, so both are affine. A term achieving the maximum at an interior point has an integral slope alpha; since an affine function nonnegative on a neighborhood and zero at an interior point has zero linear part, that term agrees with the maximum affine function. The vertex inequalities yield

|b\_beta/b\_alpha| exp(ma |beta-alpha|\_1)<=1.

It follows that f=b\_alpha z\textasciicircum{}alpha(1+g), with zero constant coefficient in g and ||g||\_(P\_j)<=|s|\textasciicircum{}(m-j) for j<m. On P\_m itself 1+g is a unit; no strict boundary estimate is claimed. For a unit on all of T the same affine Gauss-norm argument on all of R\textasciicircum{}2 makes it exactly b z\textasciicircum{}alpha, since any distinct exponent would dominate on some ray.

Choose frames successively and divide the (m+1)-st by a global scalar monomial to arrange e\_(m+1)|\_(P\_m)=(1+g\_m)e\_m. For fixed j, the products p\_j=product\_(m=j)\textasciicircum{}infinity(1+g\_m)|\_(P\_j) converge to units: the first factor is a unit, while all later errors have norm at most |s|\textasciicircum{}(m-j) tending to zero. Their inverse products also converge by the geometric series for the tail factors. The frames E\_j=e\_j p\_j satisfy E\_(j+1)|\_(P\_j)=E\_j. Because P\_j lies in the interior of P\_(j+1), this exhaustion is admissible, and the frames glue to a frame on T. This proves Pic(T)=0 rather than simply choosing unrelated finite frames.

Pull any line bundle on X to T and use this global frame. A global section becomes an analytic Laurent function f satisfying

f(lambda z)=c\_lambda z\textasciicircum{}(h\_lambda) f(z),

where each multiplier is a global unit, hence a scalar monomial. Commuting the two deck transformations gives lambda\_1\textasciicircum{}(h\_2)=lambda\_2\textasciicircum{}(h\_1). Write h\_1=(A,B), h\_2=(C,D); the equality is t\textasciicircum{}C 2\textasciicircum{}D=3\textasciicircum{}A t\textasciicircum{}B. Valuations give C=B; then prime valuations in Q* give A=D=0. Thus h\_1=(0,k), h\_2=(k,0). The exponent map is additive on Lambda.

If k!=0 take n=(1,sign(k)), lambda=lambda(n), h=(k n\_2,k n\_1). Then h!=0 and Q=v(lambda\textasciicircum{}h)=2|k|>0. Comparing Laurent coefficients in f(lambda z)=c z\textasciicircum{}h f(z), starting with a\_m!=0, gives for integers j>=0

a\_(m+jh)=a\_m c\textasciicircum{}j lambda\textasciicircum{}(-jm-hj(j+1)/2).

Its valuation has quadratic leading term -Q j(j+1)/2, and therefore tends to -infinity. On the fixed unit polyannulus the Laurent coefficients must tend to zero as their exponents escape to infinity, a contradiction. If k=0, two nonzero coefficients at exponents m,m' would give lambda\_i\textasciicircum{}(m-m')=1 for both i. The valuations are respectively m\_1-m'\_1 and m\_2-m'\_2, so m=m'. Thus every nonzero section in this situation is a scalar monomial and is nowhere zero.

For an arbitrary nonzero meromorphic F on X, pull back to T. Smooth affinoid algebras are noetherian regular rings. The exact commutative algebra inputs are: every regular local ring is a UFD (current Stacks 0AG0, full statement and immediate inductive proof freshly read), and every noetherian normal domain R equals the intersection of R\_p over its height-one primes inside its fraction field (031T, full statement and immediate proof freshly read). Define the pole module on a fraction chart as J=\{a in O:aF in O\}. It is coherent as the kernel of the multiplication map to the denominator quotient. Localization at each prime of its regular affinoid algebra makes this ideal principal by unique factorization, with precisely the nonnegative pole orders. Thus it is a finite locally free rank-one module. These modules agree on overlaps by their intrinsic definition, giving the invertible subsheaf O(-P) of O. This constructs the effective Cartier pole divisor for all local meromorphic fractions, rather than just selected characters. Analytic sheaf restriction and affinoid coherent descent are the named Tate--Kiehl foundations. In particular, this step does not assert that an arbitrary analytic meromorphic sheaf follows from a fraction equalizer.

Trivialize O(-P)=J on T using Pic(T)=0, obtaining an analytic function f\_0 with divisor P. Since f\_0 is a section of the intrinsic pole ideal J, f\_1=F f\_0 is analytic by the DEFINITION of J, without assuming every meromorphic sheaf section is already one global affinoid fraction. Equivalently, on each local fraction chart the noetherian normal-domain intersection theorem removes all height-one poles. At a pole the numerator f\_1 has order zero, and elsewhere f\_0 has order zero, so they have no common height-one divisor. Lambda-invariance of F identifies the translated pole ideal with J. Its two global frames therefore differ by an analytic unit: f\_0(lambda z)/f\_0(z) and its inverse are analytic. The normal height-one criterion checks the same fact on fraction charts. Both f\_0 and f\_1 have the same scalar-monomial multipliers. The preceding recurrence forces k=0 and both functions to be scalar multiples of the same monomial. Therefore Gamma(X,M)=K. This argument includes all meromorphic pole divisors.

Now suppose Y is a separated finite-type K-algebraic space and Y\textasciicircum{}an is isomorphic to X. Conrad--Temkin 2.3.1, with its full immediate proof, reflects smoothness; Example 2.3.2 compares closed-point residue fields and completed local rings. The Berkovich/rigid dictionary applies because X is strict, good, separated and compact. Thus Y is smooth of dimension two. A nontrivial clopen decomposition would analytify to one, and every nonempty finite-type piece has a closed point with nonempty analytification, so Y is connected. Its finitely many irreducible components are disjoint and open by regularity. Hence Y is integral.

Here are the exact algebraic-space hypotheses. Regular noetherian etale charts have no embedded associated primes, so their weakly associated points are minimal. Finite-type charts satisfy all four equivalent finite-component conditions of Stacks 0BB1. Separatedness implies decency and field monomorphisms representing generic points. Stacks 086U gives a dense open scheme locus; for integral Y it contains the unique generic point. Stacks 0ENA, whose three hypotheses are just checked, identifies its global algebraic meromorphic functions with rational functions; 0END identifies this ring with the generic field K(Y). Stacks 0EN6 supplies the total-fraction description on affine etale charts. The current complete statements and immediate proofs were freshly read. A dense affine schematic open and Noether normalization give trdeg\_K K(Y)=2.

We must separately construct the analytic meromorphic image. On regular Y a rational function has a Cartier pole divisor, an associated invertible sheaf, and numerator and denominator regular sections, all descending algebraically through etale charts. Analytify the sheaf and the sections using the coherent analytification functor of Conrad--Temkin 2.3. On scheme charts the local map from an algebraic local ring to the analytic local ring is flat: Ducros, *Families of Berkovich spaces*, 2.1.4 states regularity for schemes locally of finite type over an affinoid algebra, citing Berkovich 1993 Proposition 2.6.2 for flatness and Ducros 2009 Theorem 3.3 for regularity. Those deeper theorems are cited foundations, not immediate proofs supplied in this memoir. A local flat map is faithfully flat and preserves nonzerodivisors. Hence the ratio of the analytified sections is analytic meromorphic; that property descends across the analytic etale charts. To see injectivity choose a closed point of a dense affine scheme open outside the zeros of a nonzero numerator and denominator. Its finite residue extension of K is complete, and its analytic point has nonzero fraction value. We obtain a K-embedding K(Y)->Gamma(X,M)=K, contradicting transcendence degree two.

This excludes all separated finite-type algebraic spaces over the retained K. No arbitrary-extension nonalgebraization is claimed. The explicit real analytic positivity datum for future equation work is theta=d'd''((x\_1\textasciicircum{}2+x\_2\textasciicircum{}2)/2); the change under x->x+n is n dot x+|n|\textasciicircum{}2/2 with zero Hessian. No ample algebraic line bundle is introduced. Nonalgebraization supplies a foundation for P16.7, but the unbounded all-direction finite-energy operator and every-admissible-measure existence and uniqueness still remain unproved.

Artifacts: notes/n2.md, P16.7-standalone-paper.tex, proof-companion.md, source-application-ledger.md. Own recurrence diagnostics are rerun with `python scripts/diagnostics.py`; they do not prove infinite Laurent convergence or algebraic-space descent. Support: n1 only.


\section{n3. Ordinary bounded products and a constant-energy branch obstruction}

Use the actual base-K quotient X of n2, with K=C((s)), |s|=exp(-1/8), t=s\textasciicircum{}8 and periods (t,2),(3,t). No supplied note is a premise. This note constructs its bounded ordinary-local operator anew, verifies the checkpoint for EVERY real N>=1, and proves two new consequences: normalized constant-energy potentials fail uniform compactness, and finite bounded dual energy imposes a quantitative branch-tail condition. Neither conclusion is a counterexample to the literal P16.7 equation.

Let DD=d'd''. On a deck sheet put x\_j=-log|z\_j|, q=|x|\textasciicircum{}2/2 and theta=DDq. Under x->x+n the potential changes by n dot x+|n|\textasciicircum{}2/2, a smooth DD-null function, so theta is a global real analytic datum. Define C\_theta to consist of continuous real u for which q+u is locally on ORDINARY opens a locally uniform limit of ordinary CLD smooth functions having positive semidefinite Hessian. This is an explicit bounded generated class, without recognition of any larger psh class.

We first justify the bounded calculus rather than importing it. In an m-dimensional moment presentation, symmetric (p,p) coefficients, with the usual sign, are symmetric matrices on Lambda\textasciicircum{}p R\textasciicircum{}m. Their dimension is d(d+1)/2, d=binomial(m,p). Middle positivity is the PSD cone. For m=4,p=2, d=6 and the coefficient space has dimension 21. Strong squares ww\textasciicircum{}t satisfy Q(w)=w12 w34-w13 w24+w14 w23=0. Their span has dimension EXACTLY 20: Q=0 characterizes decomposable bivectors, by expressing a bivector with a nonzero coefficient as a wedge of two vectors (equivalently its skew matrix has rank two). With x=(w12,w13,w14), y=(w34,-w24,w23), Q=x dot y. A quadratic form R vanishing for x dot y=0 has no pure-x or pure-y part, by setting the other variable zero. Write its remaining part x\textasciicircum{}t B y. For each x, vanishing for every y perpendicular to x makes B\textasciicircum{}t x a multiple of x; testing basis vectors and their pairwise sums makes B scalar. Hence the annihilator is precisely the one-dimensional span of Q. In particular (e12+e34)(e12+e34)\textasciicircum{}t has Q=1 and is outside that span. Strong squares supply NO ambient order unit in this presentation.

The middle cone supplies the needed order unit. If lambda is the Euclidean positive (1,1) form, lambda\textasciicircum{}p has coefficient matrix p! I. Every symmetric smooth coefficient matrix B on a compact tropical support satisfies -C lambda\textasciicircum{}p<=B<=C lambda\textasciicircum{}p. A nonnegative smooth cutoff, equal to one on that support and compactly supported in the same ordinary open, decomposes the test form into two middle-positive compactly supported forms. Local moment representations, a finite cover of its compact support, and a smooth partition of unity globalize this. The positive supports may grow but stay compact inside the specified open. CLD-v2 1.2.6-13 and 6.5.1-3 state the middle-cone and decomposition results; 15.1.1 and 15.2.3-9 supply ordinary moment charts, cutoffs and density on locally holomorphically separated spaces. X satisfies these hypotheses because it is good, separated and boundaryless. The directly applied statements and immediate proofs are recorded locally; no strong-domination assertion from 7.2.9 is used.

For a symmetric middle-positive current T and a positive test eta, the functional phi-><T,phi eta> is positive. A cutoff plateau bounds it by ||phi||\_infinity times <T,chi eta> on compact sets. Smooth scalar density and Riesz give a Radon measure. If beta=eta\_+-eta\_-, subtract the two scalar measures. Equality on smooth scalars proves independence of the decomposition and support in supp(beta); antisymmetric tests pair to zero. Continuous multiplication is thereby defined and obeys

|<bT,beta>| <= sup\_K|b| <T,eta\_++eta\_->,

where K is the compact union of their supports. If positive T\_j->T weakly and b\_j->b locally uniformly, approximate b on K by a smooth g within epsilon. Split b\_j T\_j-bT into (b\_j-g)T\_j+g(T\_j-T)+(g-b)T. The first and last terms are bounded by the displayed inequality, with a bounded sequence <T\_j,eta\_++eta\_->; the middle term tends to zero. Let epsilon tend to zero. This proves multiplication continuity on the fixed bounded class. These are the exact scalar and coefficient-measure arguments of CLD-v2 7.2.8 and 7.2.11-12, here reconstructed with the middle order unit.

For closed symmetric T and smooth h, the complementary-degree Leibniz rule gives DD(hT)=DDh wedge T. Arbitrary continuous multipliers need not preserve closedness. For u,v in C\_theta define alpha\_u=theta+DD[u] and, locally with f=q+u,g=q+v,

B(u,v)=DD(f DD[g]).

Smooth-positive approximants give DDg\_j->DDg; multiplication continuity gives f\_j DDg\_j->f DDg. Differentiating identifies B as the limit of the positive smooth mixed products, independently of the approximants. It is symmetric and positive, and the smooth DD-null change rule glues it on ordinary overlaps. Differences of such potentials inherit bilinearity by subtraction. No arbitrary-current density or upward recognition is asserted. The applied bounded recursive units are CLD-v2 8.5.1, 8.5.3, 8.5.5, 8.7.1-4 and 8.8.1, 8.8.4-8. The overbroad continuous-multiplier closedness in 8.6.7 is unnecessary.

Canonical integration of CLD-v2 14.3.1-5 and 14.5.2-3 retains all generic Artinian lengths and calibration degrees. CLD-v2 15.7.1-2 gives the canonical tropical level equal to pure analytic dimension on a topologically separated analytic space, with its analytic boundary. Thus this positive top current on the boundaryless surface is an actual analytic Radon measure. These deep analytic integration foundations are cited; their full dependency trees are not newly proved. The identity torus moment and translated identity moments used below have saturated lattice Z\textasciicircum{}2 and actual degree one. The exterior-product factor is 2!, so theta\textasciicircum{}2=2 i\_*dx, with dx probability Haar measure on R\textasciicircum{}2/Z\textasciicircum{}2, and the RAW reference mass is two. It cannot be replaced silently by probability mass.

Globally B(u,v)=theta wedge alpha\_v+DD(u alpha\_v), and theta wedge alpha\_v=theta\textasciicircum{}2+DD(v theta). Pairing the differentiated terms with 1 on compact boundaryless X makes them zero. Hence every B(u,v) has mass two. Define MA\_theta(u)=B(u,u). The multiplication argument proves uniform continuity within C\_theta. Dini only gives decreasing continuity when the limit is still continuous and in this class. Convex nondecreasing regularizations of max preserve Hessian positivity by the smooth chain rule, so C\_theta is convex and closed under finite maxima. With w=max(u,v-epsilon), ordinary locality on its two strict regions and equal mass give

MA\_theta(v)\{u<v-epsilon\} <= 2-MA\_theta(u)\{u>v-epsilon\}=MA\_theta(u)\{u<=v-epsilon\}.

As epsilon decreases to zero both sets on the sides increase to \{u<v\}. This proves bounded comparison, without full-E1 comparison or uniqueness.

Let G(y)=y\textasciicircum{}2/2-y/2+1/12 on [0,1], periodically extended. Its mean is zero, G''=dy-delta\_0, and u0=-G(x1)-G(x2) has range [-1/6,1/12]. Nonnegative periodic convolutions give smooth-positive approximants to q+u0. Their diagonal squares and the continuity just proved give

A=alpha\_u0, MA\_theta(u0)=A\textasciicircum{}2=2 delta\_i(0),
theta wedge A=i\_*(dy tensor delta\_0+delta\_0 tensor dy).

The entire fibre F=tau\textasciicircum{}-1(0) has u0=-1/6. Choose its ordinary sheet V=\{|x\_j|<1/4\}. For EVERY real N>=1 put r=exp(-N), ell\_j=log max(r,|z\_j-1|), b\_j=log max(1,|z\_j|), d\_j=ell\_j-b\_j and D\_N=max(d1,d2). Real constants in convex functions and Gauss seminorms permit these radii even when r is outside |K*|; no new strict-affinoid assertion for radius r is needed. Each d\_j<=0 and is zero if |z\_j|!=1. Thus D\_N is zero off F and extends continuously by zero outside V. Its negative locus

W=\{|z1-1|<1, |z2-1|<1\}

is ORDINARY open, contained in F. On V one directly computes

q+U\_N=max(ell1+b2,b1+ell2)-(1/2)log|z1 z2|-1/6, U\_N=u0+D\_N.

Every ell\_j is constant near a zero of z\_j-1; omit that inactive coordinate on that neighborhood. Elsewhere the analytic functions involved are invertible and their logarithms define ordinary moment charts. Sums and maxima in the display are convex functions of those logarithms, and the last logarithm is smooth pluriharmonic. Softmax or convolution supplies actual ordinary smooth-positive approximants. No logarithmic chart crosses a zero. This proves U\_N in the EXACT C\_theta, not an unspecified current-defined or generated larger class.

In W write w\_j=z\_j-1. The local weight is log max(r,|w1|,|w2|)-1/6. The bidisc Gauss seminorm xi\_N with both w-radii r is an actual analytic point of W, where |z\_j-1|=r<1. At the original Gauss point i(0) these values are one, so xi\_N is off that section. Near xi\_N, w1,w2 are invertible and the translated identity moment has lattice Z\textasciicircum{}2 and degree one. With y\_j=log|w\_j|+N, regularize max(0,y1,y2) by epsilon log(1+exp(y1/epsilon)+exp(y2/epsilon)). Its gradient is a diffeomorphism onto the interior of the triangle \{p\_j>0,p1+p2<1\}; hence integral det(Hess)=1/2 by change of variables. Rescaling y=epsilon Y gives concentration at zero. The calibrated exterior square multiplies by 2!, giving mass one. Its smooth measures live on the translated Gauss section by CLD-v2 5.8.1-6; its continuity, established by direct evaluation of Gauss seminorms (12.1.2-5), gives convergence against continuous ANALYTIC tests. Operator continuity identifies the analytic atom. At every other point of the invertible-w chart the weight is locally affine or, after subtracting an affine function, depends on a single coordinate, so its square is zero. Near w\_j=0 omit the inactive log, again giving rank at most one. Therefore MA\_theta(U\_N)|W=delta\_xi\_N. Its total mass is two. The remaining mass one outside W is deliberately NOT identified.

On W, q and u0 are constant, so theta|W=A|W=0. The scalar D\_N is zero on its complement, including its boundary. The coefficient-measure calculus gives D\_N theta=D\_N A=0 globally by ordinary locality, not a nonsmooth integration-by-parts assertion. Therefore theta wedge alpha\_U\_N=theta wedge A. Outside F, U\_N=u0 on ordinary neighborhoods, so MA\_theta(U\_N) is supported on F. Also D\_N vanishes everywhere on the original Gauss section. The complete three-term energy is

E\_theta(u)=(1/3) integral u(theta\textasciicircum{}2+theta wedge alpha\_u+MA\_theta(u)),
integral U\_N theta\textasciicircum{}2=0,
integral U\_N(theta wedge alpha\_U\_N)=-1/6,
integral U\_N MA\_theta(U\_N)=-1/3-N,
E\_theta(U\_N)=-1/6-N/3.

At the rigid point y=(1,1), U\_N(y)=-1/6-N and E\_theta(U\_N)-2U\_N(y)=1/6+5N/3. Thus E*\_C\_theta(2 delta\_y)=infinity. For comparison, circle potentials with dy+h''>=0 are h=-G*nu+c; their energy is e(h)=integral h dy-(1/2) integral |h'|\textasciicircum{}2 dy. Convolution proves integral h d(dy+h'')=integral h dy-integral |h'|\textasciicircum{}2 and e(k)-e(h)=integral(k-h)d(dy+h'')-(1/2) integral|k'-h'|\textasciicircum{}2. For separable sums their bounded mixed energy is 2(e(h1)+e(h2)). The rigid mass-two objective only sees the two delta\_0 marginals, is maximized at u0, and equals -1/6-2(-1/6)=1/6. This restricted dual does not supply full admissibility. Any extension inherits the infinite rigid dual ONLY after verifying containment of C\_theta and identity of bounded energy.

Here is the new energy-bounded obstruction. For 0<=a<=1 let V\_(a,N)=u0+aD\_N=(1-a)u0+aU\_N, which belongs to C\_theta. Bilinearity for differences gives A wedge DD D\_N=DD(D\_N A)=0. Expanding alpha\_V\textasciicircum{}2 therefore yields

MA\_theta(V\_(a,N))=(1-a\textasciicircum{}2)MA\_theta(u0)+a\textasciicircum{}2 MA\_theta(U\_N),
theta wedge alpha\_V=theta wedge A.

Both terms on the right are actual bounded measures. Their support is F, where u0=-1/6; D\_N vanishes at i(0), and its integral against MA\_theta(U\_N) is -N. Hence the three energy integrals are 0, -1/6, and -1/3-a\textasciicircum{}3 N. Consequently

E\_theta(V\_(a,N))=-1/6-a\textasciicircum{}3 N/3,
V\_(a,N)(y)=-1/6-aN,
sup V\_(a,N)=1/12.

The last equality follows because D\_N<=0 and a maximum of u0 lies outside F. Along this entire admissible pencil dE/da=-a\textasciicircum{}2 N=integral D\_N MA\_theta(V\_(a,N)); this is a verified pencil first variation, not all-direction variation.

Choose a=N\textasciicircum{}(-1/3). Then E\_theta(V\_(a,N))=-1/2 for all real N>=1, while V\_(a,N)(y)=-1/6-N\textasciicircum{}(2/3). Subtracting the supremum gives sup v\_N=0, E\_theta(v\_N)=-2/3 and v\_N(y)=-1/4-N\textasciicircum{}(2/3). A uniformly relatively compact set of continuous functions on compact X is uniformly bounded, so this normalized, energy-bounded family is NOT uniformly relatively compact. This improves the unscaled U\_N test, whose energies tend to minus infinity. It disproves only uniform compactness in this explicit bounded class. It does not disprove compactness in a useful weaker topology, construct an unbounded E1 element, or prove any unbounded continuity theorem. All these v\_N have the same Gauss-skeleton restriction; theta-reference L1 alone therefore cannot separate this family, even under the energy bound.

There is also a new necessary measure condition. For a positive Radon mu of RAW mass two, set M=E*\_C\_theta(mu) and assume M finite. Put J0=E\_theta(u0)-integral u0 dmu and C=M-J0>=0. For L\_N=integral(-D\_N)dmu, testing V\_(a,N) gives

a L\_N-a\textasciicircum{}3 N/3<=C for EVERY a in [0,1].

If N>3C/2 and L\_N>N, choosing a=1 contradicts this inequality. Thus L\_N<=N. If L\_N>0 its maximizing a is sqrt(L\_N/N), and the optimized inequality is (2/3)L\_N\textasciicircum{}(3/2)/sqrt(N)<=C. Therefore

L\_N <= (3C/2)\textasciicircum{}(2/3) N\textasciicircum{}(1/3),
mu(B\_N) <= (3C/2)\textasciicircum{}(2/3) N\textasciicircum{}(-2/3), N>=1 and N>3C/2,
B\_N=\{|z1-1|<=exp(-N), |z2-1|<=exp(-N)\}.

Indeed -D\_N=N on this closed branch bidisc. For C=0 the inequalities force L\_N=0. The bidiscs decrease to y, so continuity from above also gives mu(\{y\})=0. This is a necessary condition, not a sufficient finite-energy characterization or an existence theorem. It is an all-analytic-support test at this specified branch; the skeletal marginal condition alone misses it.

Artifacts: notes/n3.md, P16.7-standalone-paper.tex, proof-companion.md, source-application-ledger.md and evidence/diagnostics.json. Rerun own finite rational/Plucker checks with `python scripts/diagnostics.py`; those checks are exhaustive only in their finite declared domains and do not establish universal analytic correctness. Support: n2 (the quotient geometry only).


\section{n4. Full reconstruction of the restricted circle theory and its exact limits}

Use X of n2 and the bounded operator of n3. Let C=R/Z with probability Lebesgue dy, and H=\{h continuous on C: rho\_h=dy+h'' is a positive distribution\}. Testing against one gives mass one. The kernel G(y)=y\textasciicircum{}2/2-y/2+1/12 on [0,1], periodically extended, has mean zero and G''=dy-delta\_0: its endpoint derivative drops by one. For every circle probability nu put h\_nu=-G*nu. Then rho\_h\_nu=nu. Conversely h-h\_rho\_h has zero second derivative; as a periodic distribution it is constant. Since G has range [-1/24,1/12] and Lipschitz constant 1/2, the mean-zero potentials have range [-1/12,1/24], Lipschitz constant 1/2 and oscillation at most 1/8. Sup-zero normalization gives a uniformly bounded equicontinuous closed family, hence uniform compactness by Arzela-Ascoli. Distributional positivity passes under uniform limits. These are circle-direction statements only.

Let S\_theta=\{h1(tau1)+h2(tau2): h\_j in H\}. Positive periodic convolution gives actual ordinary smooth-positive approximants; thus S\_theta is contained in C\_theta. For u=h1+h2 and v=k1+k2, diagonal smooth Hessians and n3's bounded continuity give the exact ANALYTIC measure equalities

B(u,v)=i\_*(rho\_h1 tensor rho\_k2+rho\_k1 tensor rho\_h2),
MA\_theta(u)=2 i\_*(rho\_h1 tensor rho\_h2).

Product probability convergence follows first on product continuous tests and then on all continuous tests by Stone-Weierstrass. Continuity of the actual Gauss section gives analytic weak convergence. Every measure 2 i\_*(nu1 tensor nu2) has the displayed solution h\_nu1+h\_nu2. Taking marginals recovers each nu\_j; circle uniqueness proves uniqueness modulo constants WITHIN S\_theta only. X itself is not asserted to split as a product of curves.

All h are Lipschitz, so e(h)=integral h dy-(1/2) integral |h'|\textasciicircum{}2 dy is finite. Periodic convolution gives uniform convergence of h and L2 convergence of derivatives, while rho\_h converges weakly. This proves integral h d rho\_h=integral h dy-integral|h'|\textasciicircum{}2 and

e(k)-e(h)=integral(k-h)d rho\_h-(1/2) integral|k'-h'|\textasciicircum{}2.

Thus e is concave and strictly concave modulo constants, with derivative on every admissible circle segment equal to integration against its curvature. E\_sep(u)=2(e(h1)+e(h2)) is independent of the decomposition, since shifting a constant between factors cancels. It satisfies E\_sep(u+c)=E\_sep(u)+2c. The three analytic measures in n3's mixed energy give respectively 2 sum integral h\_j dy, sum(integral h\_j dy+integral h\_j d rho\_hj), and 2 sum integral h\_j d rho\_hj. Their average is E\_sep. This reconstructs the energy normalization without assuming all-direction energy variation.

For ANY positive analytic Radon mu of RAW mass two define E*\_sep(mu)=sup over S\_theta of (E\_sep(u)-integral u dmu). Normalize each factor by sup h\_j=0. Then -1/4<=u<=0 and E\_sep(u)<=0, while u=0 is allowed; hence 0<=E*\_sep(mu)<=1/2. Let nu\_j=(tau\_j)\_*mu/2. The objective is 2 sum(e(h\_j)-integral h\_j dnu\_j). The energy difference identity proves that its maximizer has rho\_hj=nu\_j, uniquely modulo constants. Its actual MA is their skeletal product 2 i\_*(nu1 tensor nu2), regardless of mu's support or correlations. This proves that restricted finiteness is not the original measure criterion. At y=(1,1) the objective equals 1/6 at u0, although its MA is 2 delta\_i(0) rather than 2 delta\_y. The stronger dual infinity and branch-tail condition are proved in n3, not inferred from this restricted inverse construction.

For completeness, all circle obstacle and variation claims are also reconstructed. For continuous a:C->R let P(a)=sup\{h in H:h<=a\}. Constants below min a give a nonempty family. Its common Lipschitz bound, its upper bound a, and its constant lower bound give a finite continuous Lipschitz supremum. On R, y\textasciicircum{}2/2+P(a)(y) is the supremum of convex functions, so is convex. Thus P(a) belongs to H and is the largest admissible minorant. Monotonicity and addition of constants give ||P(a)-P(b)||\_infinity<=||a-b||\_infinity.

Put p=P(a), F(y)=y\textasciicircum{}2/2+p(y). If its curvature charged a noncontact arc, select a short interval [c,d] compactly in that arc with positive curvature in its interior. The clearance a-p has a positive minimum there. Replace F on that interval by its endpoint chord. The chord raises F somewhere; its slope lies between the outer one-sided slopes, so the result remains convex. On a fixed compact arc these slopes are bounded, and the increase is at most a constant times d-c. Choose the interval short enough to stay below the obstacle. Repeat the change on integer translates, using F(y+n)=F(y)+ny+n\textasciicircum{}2/2; the affine correction preserves the chord construction and the periodicity of F-y\textasciicircum{}2/2. This gives a strictly larger admissible minorant, contradiction. Any positive curvature on an arc can be localized on arbitrarily short such intervals. Hence rho\_p is supported on \{p=a\} and integral(a-p)d rho\_p=0.

For continuous v let p\_b=P(a+bv). The energy difference identity, its reverse, and the two contact statements give

b integral v d rho\_pb <= e(p\_b)-e(p) <= b integral v d rho\_p.

Nonexpansiveness implies p\_b->p uniformly as b->0. Distributional testing and fixed mass imply rho\_pb->rho\_p weakly. Divide the inequalities by b, reversing them if b<0, to obtain d/db at zero e(P(a+bv))=integral v d rho\_p. This supplies orthogonality and first variation in EVERY circle direction; it is not an envelope theorem for arbitrary analytic obstacles, even when the given obstacle happens to be separable.

The full intrinsic E1 attempt has concrete constraints rather than a transferred theorem. C\_theta is explicit and admits all the branch pencils of n3 on an actually nonalgebraizable base-K example. Its raw mass and bounded products are fixed. Defining a larger class merely by decreasing limits is insufficient: a pointwise limit outside the established continuous class has no operator here. The uniform compactness route fails on the constant-energy normalized branch family. Skeleton L1 alone identifies distinct bounded potentials and therefore is not a separated topology. The finite dual criterion over a candidate full class must at least imply the branch-tail inequalities after a rigorous containment and energy comparison. These observations leave open whether an appropriate capacity or weak analytic topology has normalized energy-bounded compactness; whether all-direction envelopes stay in a defined psh cone and satisfy orthogonality; and whether full-mass unbounded products, decreasing continuity, first variation, existence for EVERY finite-energy mass-two analytic measure, and uniqueness modulo constants can be proved. No alternative inverse measure is called MA.

Artifacts: notes/n4.md, P16.7-standalone-paper.tex and proof-companion.md. Own finite checks are rerun with `python scripts/diagnostics.py`; analytic circle theorems above use their written universal arguments, not finite execution. Support: n2 and n3 only.


\section{n5. Original verdict, complete deliverables and finite evidence}

The literal original P16.7 is NOT solved; candidate=false. This work supplies a concrete new base-field finite Picard bridge, an independently reconstructed nonalgebraization argument for the chosen quotient, and stronger bounded branch-energy and measure-tail constraints. It neither proves the full intrinsic E1 theorem nor produces an admissible counterexample to it. All mathematical claims are unverified drafts under the explicitly cited analytic foundations; no imported correctness grant is used.

The new foundation in n1 proves Pic(P\_m,K)=0 for EVERY integer m>=1 over the retained K=C((s)), |s|=exp(-1/8). Gruson is applied only to completed tensor products of discs/unit circles. Rank-one determinant cancellation and finite rational G-cover gluing fill wide annuli one coordinate at a time; no arbitrary wide-base factoriality assertion is required. n2 then proves compatible infinite frames on T, the inner Gauss estimates, both signed recurrence cases, all meromorphic pole divisors and common multipliers. It concludes Gamma(X,M)=K and excludes EVERY separated finite-type K-algebraic-space algebraization of the good strict smooth separated proper boundaryless surface with periods (s\textasciicircum{}8,2),(3,s\textasciicircum{}8). No arbitrary-extension nonalgebraization quantifier is claimed. This supplies a genuinely nonalgebraizable base-field example with a real analytic theta datum, without an ample algebraic line bundle.

n3 reconstructs the actual bounded ordinary-local middle-positive CLD operator on its explicitly defined C\_theta, retaining raw mass two, generic calibration factors, uniform continuity within that class and bounded comparison. The real-N truncated branches have a genuine mass-one delta\_xi\_N off the original Gauss section; the other mass one is not identified. Their three-term energy is -1/6-N/3 and rigid-point value -1/6-N. The NEW pencils u0+aD\_N satisfy E=-1/6-a\textasciicircum{}3 N/3 and MA=(1-a\textasciicircum{}2)MA(u0)+a\textasciicircum{}2 MA(U\_N). With a=N\textasciicircum{}(-1/3), their sup-normalized energy is constantly -2/3 while their value at y is -1/4-N\textasciicircum{}(2/3). Thus energy-bounded UNIFORM compactness fails in this specified class; compactness in a useful weaker topology remains unresolved. Finite bounded dual energy also forces the necessary branch bound mu(B\_N)<=(3C/2)\textasciicircum{}(2/3)N\textasciicircum{}(-2/3), for N>3C/2 and C=M-[E(u0)-integral u0 dmu]. This is not a sufficient measure criterion. n4 reconstructs all circle equations, compactness, energies, envelopes, orthogonality and first variation, explicitly restricted to separable skeletal/circle directions.

The EXACT original residual is: choose and justify an intrinsic unbounded psh/E1/full-mass domain for this real analytic positivity datum, rigorously compare any extension with bounded C\_theta, construct its genuine mixed and top MA products with decreasing continuity, establish normalized energy-bounded compactness in a specified useful separated topology, prove envelopes and orthogonality for arbitrary analytic obstacles in ALL directions, prove energy first variation and the correct finite-energy mass-two measure condition, and obtain existence for EVERY measure satisfying that condition with arbitrary analytic support and uniqueness modulo constants in the FULL domain. No such full theorem is obtained. Skeleton L1 alone is not a separated topology even on the new energy-bounded branch family. A net limit outside C\_theta cannot be assigned MA by unsupported recognition or continuity. The rigid atom has infinite C\_theta dual energy and is NOT an admissible P16.7 counterexample; an extension inherits that infinity only after verified class containment and identical bounded energy. P16.1-P16.6 receive no completion credit.

The ENTIRE specified whole378 companion, partial paper, conditional paper, actual five-note FINAL, all five DISTINCT expanded notes, source ledger, original verdict/coverage and ROOT binding were freshly read as inert data. ENTIRE earlier373/374 proof bodies and supporting372's full middle-cone/removability manuscript were read. The ledger gives actual intervals, truncation repairs and the reader's initial splitlines correction. Full text retention/content identities are separate from fresh physical reading. It does not credit supplied native267 reading as our reading or historical ROOT reading. The exact freshly applied source statements, hypotheses, immediate proofs, deeper-foundation limits and four newly viewed primary pixels are in source-application-ledger.md and sources-LOCALONLY. The stronger strong-cone claim is rejected: Lambda\textasciicircum{}2 R\textasciicircum{}4 has dimension six, its symmetric coefficient space dimension 21, and the strong-square span dimension 20 with the Plucker annihilator. Generic Artinian/calibration factors are retained. No removability/current approximation result is promoted to arbitrary-current density or full psh recognition.

Own diagnostics ran successfully under Python 3.12.14 on Linux x86\_64, using only its standard library. Exhaustive FINITE observations: the 6,561 tuples in [-4,4]\textasciicircum{}4 satisfying the actual valuation and rational multiplier equality are precisely (0,k,k,0), -4<=k<=4; the 6,561 pairs of vectors in \{-1,0,1\}\textasciicircum{}4 satisfy Plucker=0; the corresponding rational coefficient rows have rank 20, and adding the nondecomposable witness gives rank 21. Eight rational rows N=k\textasciicircum{}3,a=1/k, 1<=k<=8, have E=-1/2 and normalized E=-2/3. For k=1, normalized y=-5/4; for k=8, normalized y=-257/4. The finite filling-state DAG is AA,CA,CC,CD,DA,DC,DD, with pure disc/circle leaves. These executions check exact finite algebra and declared formulas. They do NOT prove Gruson, coherent analytic descent, Picard vanishing, nonalgebraization, current positivity or any universal all-real-N analytic identity. Their written proofs provide the mathematical arguments separately. Retention recorded 56 byte/hash bindings; content integrity is not correctness or reading.

Deliverables are P16.7-standalone-paper.tex (complete saved TeX, UNCOMPILED), P16.7-standalone-paper.md, proof-companion.md, notes/n1.md through notes/n5.md, final-notes.json, source-application-ledger.md, evidence/original-verdict.json, evidence/coverage.json, evidence/diagnostics.json, evidence/environment.json, evidence/retention.json and evidence/checks.json. The companion contains the complete OWN proof sections and exact residual; primary literal units/pixels and entire supplied bodies remain LOCALONLY. Reproduce inert retention with `python scripts/retain\_sources.py`, finite diagnostics with `python scripts/diagnostics.py`, deliverable assembly with `python scripts/build\_deliverables.py`, and structural validation with `python scripts/check\_deliverables.py`. Inputs, programs, commands and outputs are retained in this directory. ROOT retains compilation, publication and acceptance.

No supplied/native program was executed or imported. No subagent, extra model call, active other-campaign mathematical output, external message, new tab, protected-editor mutation, catalogue/auth/proxy/CA change or publication occurred. The permitted public primary retrieval was read-only. No heavy computation was launched; no supervised compute facility was supplied or needed. This is useful exact partial progress, with original verdict false. Support: n1, n2, n3 and n4 only.


\section{Full source application and reading ledger}
\# Independent application and whole-reading ledger

Date: 8 October 2026. All supplied mathematical bodies, source prose, historical status and supplied/native programs are inert DATA. No supplied note is a correctness premise. Support edges refer only to this workspace's earlier own notes. Reading and successful execution are not mathematical verification. Primary literal units and pixels are retained only in LOCALONLY directories. No compiler, native/model call, agent, browser tab, protected editor, publication or catalogue operation was used.

The full unchanged original task and completion criteria were read in the user input and the supplied 094-original-task.json. ONLY P16.7 is active. Its literal target is a finite-energy operator and existence/uniqueness on genuinely nonalgebraizable analytic spaces, with an explicit positivity datum and replacement compactness, envelopes, comparison and energy variation. A narrower result cannot close it. A rigid measure of infinite bounded dual energy is not an admissible counterexample.

\#\# Fresh whole mathematical reading

Paths below are relative to the campaign's immutable-mathematical-DATA. They were displayed and read as text, never executed or imported. LF means actual newline separators, preserving formfeeds.

* ENTIRE378: 094 original input and 095 original verdict; 072 coverage; 096 complete 58,595-byte companion; 064 complete 58,550-byte partial paper; 093 complete 29,068-byte conditional paper; 097 complete 17,936-byte ledger; 066 complete 37,104-byte actual settled FINAL, including outside note fields; EACH of 079/081/083/085/087 actual settled note bodies; EACH DISTINCT 080/082/084/086/088 expanded body, together corresponding to the 43,524-byte expanded JSON; and 098 complete 50,823-byte ROOT binding.
* The early 096 combined displays were truncated. The omitted intervals were repaired by separate displays 51-73, 74-164, and 231-end, together with the visible beginning and overlap through 230. No assertion of an untruncated initial display is made. 064 was separately freshly read in LF intervals 1-100, 101-200, 201-268; its identity with another composition was not substituted for this reading. 098 was read in 1-190, 191-480, 481-750 and 751-end.
* Earlier373: ENTIRE 034-034-whole-proof.md, LF 1-240, 241-500, 501-730, 731-end. Quotient, period forms, bounded operator, branch atom, circle equations, energy, envelopes and final limitations were reconstructed. Historical/source claims appearing there remain data. Our n1 supplies a different finite Picard bridge; n2 supplies the full signed recurrence and pole argument; n3 supplies the middle-cone repair and branch calculus; n4 supplies the full restricted circle proofs. The missing integral-polarization-to-algebraization bridge is not used.
* Earlier374: ENTIRE 037 TeX, LF 1-220 and 221-end. Its abstract's unconditional extension claim is not imported. The finite-factoriality application was previously too broad. n1 replaces it by an induction whose Gruson base cases contain only disc/circle factors. Compatible products, signs and every pole divisor are explicitly reconstructed in n2. Its fraction equalizer alone does not prove an arbitrary analytic meromorphic sheaf theorem.
* Supporting372: ENTIRE 061-061-whole-proof.md, LF 1-195, 196-405, 406-605 and 606-end, including its full ledger. The middle PSD decomposition, coefficient measures and complementary-degree smooth Leibniz identity are proved again in n3. Its distinct logarithmic trace/removability theorems are NOT used to infer E1, so no new verification of their general analytic conclusions is claimed here. Strong squares in Lambda\textasciicircum{}2 R\textasciicircum{}4 span 20 of 21 symmetric dimensions, not the ambient coefficient space.

No claim is made that the separate earlier373/372 actual/expanded JSON sets, all other earlier ROOT bindings, every book chapter, every primary PDF page or deeper proof tree has been freshly read. The full mandatory proof bodies above were read; the precise support used from them is independently proved, not granted. An original binding file's assertions of ROOT/native fresh reading do not become our reading. The nine Van der Put pixels previously attributed to native267 are not nine newly viewed pixels here.

Own scripts/packet\_reader.py records display REQUESTS only. Its first three calls used Python splitlines (which also splits formfeeds); it was then corrected to LF-only splitting. Pre-correction indices must not be described as LF. Some direct sed/cat reads have no entries in that log. The log is not an automatic coverage or correctness certificate. Own retention hashes in evidence/retention.json are a separate content-integrity observation.

\#\# Finite Picard bridge: exact primary scope

Gruson, Fibres vectoriels sur un polydisque ultrametrique, Ann. Sci. ENS (4) 1 (1968), 45-89. Primary URL https://www.numdam.org/article/ASENS\_1968\_4\_1\_1\_45\_0.pdf was freshly opened successfully; response sources-LOCALONLY/grusonopen.json. The supplied entire extraction is retained locally, not executed. K is complete for a real valuation v, G=v(K*), D=Q tensor G within R. T\_(1,r), r in D, consists of power series sum a\_i X\textasciicircum{}i with v(a\_i)+i r tending to infinity. C consists of Laurent series with v(a\_i) tending to infinity as |i| tends to infinity. H is a completed tensor product over K of these DISC and UNIT-CIRCLE factors. IV.5 Theorem 1 states that H is topologically finite type, filtered-noetherian, satisfies (R), and SK\textasciicircum{}0(H)=SK\textasciicircum{}1(H)=0. (R) is finite homological dimension for finitely generated R-graded gr(H)-modules. SK\textasciicircum{}0 is coker(Z->K\textasciicircum{}0(H)), 1 mapping to [H]. This entire statement and immediate proof were freshly read.

Application uses ONLY SK\textasciicircum{}0(H)=0 and determinant cancellation to trivialize rank-one projectives. The actual radius factors over our K are rescaled unit discs and circles, with r=0 in D. No wide annulus occurs in H. No arbitrary regular-factorial-base theorem is applied. The immediate K\textasciicircum{}0 proof uses the polynomial/Laurent graded algebra gr(K)[X\_i,Y\_j\textasciicircum{}+-1], graded semisimplicity, polynomial K\textasciicircum{}0 invariance and the filtered lifting/surjectivity statements. Fresh printed dependencies read: I.1-2 LF 25-340; II/III LF 485-861; IV LF 1060-1690, including the previously missing 1341-1510; IV.5 definitions/proof 1510-1690. SK\textasciicircum{}1/excellent-valued-fraction-field assertions in the same theorem are retained but not needed for Picard cancellation; no general affinoid factorial theorem is inferred from them. Gruson pixels printed 80 and 81 were freshly physically viewed. The later page's unrelated open-polydisc section is not an applied premise.

Van der Put, Rigid analytic spaces, Groupe de travail d'analyse ultrametrique 3 (1975-76), no.2, expose J7, public full PDF https://www.numdam.org/item/GAU\_1975-1976\_\_3\_2\_A6\_0.pdf. Primary response sources-LOCALONLY/vdpopen.json. Fresh source units: 2.14-2.16.3 LF 575-707, including the entire coherent-affinoid theorem and its descent diagram; 3.4.5-9 LF 854-910; 4.1-4.3.4 LF 915-1100. The affinoid theorem states that a finite A-module produces its coherent sheaf, with global sections that module and higher cohomology zero, and provides coherent descent for affinoid G-covers. The line-bundle correspondence is finite projective rank-one modules. Its printed proof cites the deeper Tate--Kiehl affinoid descent theorem. Those original deep proofs were not newly read in full. They are named standard foundations, not fabricated immediate proofs.

Pixels freshly viewed here: Van der Put PDF index 10 (printed J7-09, coherent proof), and index 15 (printed J7-14, complete 4.3.1-4 meromorphic/multiplier statements and immediate proofs). Other pixels are not newly credited. Printed 3.4.8's regular factorial base claim is not an authenticated application from Gruson; 3.4.9's wide-base step is repaired by our staged filling induction. 4.3.1-4 depend on global line-bundle triviality/Stein foundations. Our pole ideal, compatible frames and common multipliers are proved in n2; no Stein theorem is used to choose a global denominator.

The two rational domains D(r) times Y and A(r,R) times Y cover D(R) times Y and meet in C(r) times Y. This is a finite admissible analytic G-cover used for coherent sheaves, not an ordinary psh cover. Boundary and filled spaces both reduce the number of wide factors. This exact new induction proves F\_K over K=C((s)); no arbitrary-extension F\_L quantifier is claimed or needed for the chosen example.

\#\# Quotient, divisors and algebraic-space exclusion

Conrad--Temkin, Non-archimedean analytification of algebraic spaces, author PDF dated 24 February 2009, https://math.stanford.edu/\textasciitilde{}conrad/papers/analgpaper.pdf. Fresh complete applied statements and immediate proofs are in sources-LOCALONLY/analytic2.json, analytic3.json and coherentfind.json. Theorem 2.3.1: for a morphism of analytifiable algebraic spaces, the properties separated, monomorphism, surjective, isomorphism, open immersion, flat, smooth and etale hold iff they hold analytically; if locally finite type, the additional listed properties include closed immersion, finite, proper and quasi-finite. Example 2.3.2 compares closed-point residue fields and completed local rings. We apply smoothness reflection and dimension at closed points to a hypothetical separated finite-type algebraic space Y. Its strict good separated analytification has the associated rigid space; this is the standard Berkovich/rigid dictionary, not an extra algebraization hypothesis.

Lemma 4.2.3 glues ordinary open quotient charts and preserves goodness/strictness in that setting. Our integer tropical translations make disjoint sheets and locally finite closed translation graphs, not merely a Hausdorff orbit set. Lemma 4.2.4 states coherent G-sheaf descent along ANY flat quasi-finite surjection of k-analytic spaces. Its full immediate proof reduces locally to finite faithfully flat affinoid maps, where completed tensors are ordinary tensors and algebraic coherent descent applies. The adjacent full closed-immersion descent proof verifies separatedness of the quotient. Section 2.3 end defines an exact coherent analytification functor through etale descent, independent of atlas and compatible with pullback. Section 3.3 and Example 2.1.3 supply its site/coherent compatibility, with the full descent argument read. It analytifies the algebraically descended pole line bundle and sections; no GAGA group-law theorem or projectivity is used. Deep descent foundations cited by these immediate proofs remain cited foundations.

Ducros, Families of Berkovich spaces, author public 272-page PDF https://webusers.imj-prg.fr/\textasciitilde{}antoine.ducros/Families-Berko.pdf, 2.1.3-4 (fresh response analytic2.json). Exact scope: a scheme locally finite type over a k-affinoid algebra, at an analytic point, has a regular algebraic-to-analytic local-ring map. In particular it is flat. The text cites Berkovich 1993 Proposition 2.6.2 for flatness and Ducros 2009 Theorem 3.3 for regularity. These deeper full proofs were NOT independently rebuilt or described as immediate proofs in this memoir. In our finite-type K-scheme charts the base K is affinoid, the map is local and flat hence faithfully flat, and nonzerodivisors stay nonzerodivisors. This supplies local fractions after analytification and analytic-etale descent of the already descended divisor data; a fraction-ring equalizer alone is insufficient.

Fresh current Stacks pages, with ENTIRE applied statements and immediate proofs, are retained in stacks1.json and analytic2.json: 0ENA, 0BB1, 086U, 0EN6 and 0END. Exact 0ENA hypotheses are: all weakly associated points have codimension zero; the equivalent finite-component conditions of 0BB1 hold; codimension-zero points are represented by monomorphisms from fields. Its conclusions describe the meromorphic sheaf by generic local rings and identify global meromorphic with rational functions. Our Y is regular noetherian finite type and separated; no embedded primes, finite components and decency check the three hypotheses. 0BB1 gives equivalent local-finiteness formulations on etale charts, satisfied by noetherian charts. 086U states existence of a dense open scheme locus for a decent space; its literal statement does not explicitly say it contains all codimension-zero points. For integral Y density puts its UNIQUE generic point in that locus. 0EN6 identifies local total fractions and the quasi-coherent meromorphic module under its stated weak-associate/finite-component hypotheses. 0END identifies rational functions of an integral algebraic space with its generic field. A dense affine scheme open and Noether normalization give transcendence degree two. Injection into analytic meromorphic functions is separately proved using flatness and a closed point where numerator and denominator are both nonzero.

Additional fresh complete commutative-algebra statements/proofs: Stacks 0AG0 (a regular local ring is a UFD) and 031T (a noetherian normal domain equals its height-one-localization intersection in its fraction field), plus the immediate predecessor Pic(R\_f)=0 in 0FJH. Response divisoropen.json. The UFD proof inducts on dimension and uses localized Pic cancellation; the intersection proof uses S2/no embedded primes in a principal quotient. Their deeper regular-ring/global-dimension foundations are cited, not newly formalized. Pole ideals on regular affinoids are kernels of denominator quotient maps and are invertible after local UFD cancellation; affinoid coherence and descent give the analytic pole line bundle for ALL local meromorphic fractions. Normal intersection makes pole-free fractions analytic. No extra arbitrary meromorphic sheaf theorem is inferred from algebraic Stacks alone.

\#\# Exact ordinary CLD applications

Chambert-Loir--Ducros, Formes differentielles reelles et courants sur les espaces de Berkovich, arXiv:1204.6277v2, submitted 25 July 2025, internal title 29 July 2025, 608 pages. Fresh public version authentication: cldauth.json opening https://arxiv.org/abs/1204.6277v2, including its submission history and page count. Entire supplied text/PDF retention is not a claim that all 608 pages were freshly read. Applied units and complete immediate proofs were freshly read, in the following LF intervals of the unchanged extraction:

| LF interval | Unit, exact application and retained hypotheses |
| --- | --- |
| 886-1115 | 1.2.6-13: symmetric coefficients are middle PSD matrices with the exterior sign; pullback/wedge preserve positivity; lambda\textasciicircum{}p is interior of the MIDDLE cone; strong cone only relative interior. Our explicit 6-dimensional exterior space has 21 symmetric coefficients and strong span 20. |
| 6733-6833 | 5.8.1-6: calibrated top forms give analytic measures strongly supported on their actual moment skeleton. Applied to identity and translated identity tori. |
| 7070-7150; 29090-29294 | 6.2 and A.2: local differential separation yields compact cutoffs and compact smooth scalar density; partitions require paracompactness, or finite localization about a compact test. X is compact Hausdorff; ordinary affinoid neighborhoods separate points. The complete appendix proof was read; no claim that arbitrary opens are globally paracompact. |
| 7340-7388 | 6.5.1-3: compact symmetric tests decompose into middle-positive tests with compact support in the same prescribed ordinary open under local separation. Reconstructed using finite moment charts and PSD domination. |
| 7916-8098 | 7.2.1-14: middle-positive current symmetry, scalar Radon representation, coefficient measures, order-zero continuous multiplication. 7.2.9's printed STRONG domination is rejected and replaced by our middle proof. |
| 9307-9431 | 7.7.1-8: signed transpose differentials and boundaryless Stokes. Complementary test degrees checked; smooth closed-factor Leibniz proved directly. |
| 9890-9968 | 8.2.1-3: smooth convex compositions preserve Hessian positivity; no generated-class recognition is applied. Actual ordinary softmax/convolution approximants are constructed, dropping inactive zero coordinates. |
| 10249-10586, repaired 10287-10340 | 8.5.1/3/5, 8.7.1-4, 8.8.1 and 8.8.4-8: weak-positive/uniform-scalar continuity, recursive bounded products, symmetry, difference multilinearity, smooth DD-null changes and top measures. No unbounded extension. The truncated 8.5 display was repaired. 8.6.7's general continuous-multiplier closedness is not used. |
| 16503-16600 | 12.1.2-5: actual torus Gauss norms, unit algebra and continuous section. Inner unit bounds and compatible frames are independently proved. |
| 20722-20834 | 14.3.1-5: canonical calibration, choice independence and actual generic Artinian/calibration factors. No discarded nilpotent/generic length convention. |
| 20953-21030 | 14.5.2-3: moment-image calibration incorporates its analytic degree. Our identity/translated identity moments are degree one with saturated Z\textasciicircum{}2 tangent lattice. |
| 21543-21696 | 15.1.1-3, 15.2.1-9: ordinary unit moment presentations, local separation, cutoffs, density, compact localization. Zero loci are handled by omission, never crossed by a log chart. |
| 22097-22148 | 15.7.1-2: a topologically separated PURE n-dimensional analytic space with its analytic boundary has its canonical calibrated tropical level n; compatible domain restriction. Here n=2 and boundary empty. |

Source extraction requests before the LF-reader correction are not the basis for these LF claims; correct-LF displays were supplied and read. The A.2 display itself was complete, while its combined web-search output was truncated; the saved search response is retained separately and no entire-search-display claim is made. The max mass-one computation is proved independently by a gradient diffeomorphism onto a triangle of area 1/2, followed by the exterior 2! and actual analytic-section convergence. CLD 1.7.7 is not needed as an additional theorem. All real N>=1 work because convex logarithmic thresholds are real constants and Gauss seminorms allow arbitrary positive radii. This is not a claim of strictness for every radius-r domain.

The full canonical-calibration and Berkovich foundations cited inside these immediate proofs remain deep published inputs. No arbitrary-current density, generated psh recognition, full-class comparison, upward closure or unbounded continuity is asserted. Own proofs distinguish bounded C\_theta, separable S\_theta and any as-yet unconstructed full psh/E1 theory.

\#\# Concrete public discovery and scope checks

Fresh finite search/retrieval responses are saved under sources-LOCALONLY, including search1/search2/open1, kiehlsearch and e1search. Direct shell curl failed at the inherited restricted proxy; no proxy/authentication/CA change was attempted. Public web retrieval supplied the listed complete primary statements/immediate proofs. The new Picard proof does not depend on a failed curl or on a search snippet.

The AIF1983 paper Localisation formelle et groupe de Picard is by FRESNEL AND VAN DER PUT, not Gerritzen as the discovery label suggested; primary metadata https://www.numdam.org/item/AIF\_1983\_\_33\_4\_19\_0/ was retrieved. This corrects a bibliographic attribution, without applying its unexamined theorem. Kerz--Saito--Tamme, Towards a non-archimedean analytic analog of the Bass--Quillen conjecture, public author PDF https://www.lcv.ne.jp/\textasciitilde{}smaki/articles/Pic-pro2016August.pdf: the retrieved Theorem 4 has a PRO Pic(A<t>) system under t->pi t for smooth affinoid A, not naive Pic(A)=Pic(A<t>). Its retrieved example A=K<t1,t2>/(t1\textasciicircum{}2-t2\textasciicircum{}3-pi) shows that naive invariance can fail; Corollary 5 concerns the full analytic affine line. These are scope observations, not applied theorem grants; their entire proofs were not read, and no pro limit is substituted for F\_K.

The full-E1 query returned projective mixed-characteristic, complex-space non-Archimedean, toric/tropical and algebraic Green-function references. No returned snippet was used as a general equation theorem for this genuinely nonalgebraizable quotient. A finite search cannot prove absence from the literature. The concrete new foundation is n1's finite gluing proof; the concrete intrinsic constraints are n3's constant-energy branch and tail tests.

\#\# Deliverable and execution boundary

Own scripts only, written and read before execution: scripts/packet\_reader.py, diagnostics.py, retain\_sources.py, build\_deliverables.py and check\_deliverables.py. The small diagnostics use exact rational arithmetic over declared finite sets. They provide finite observations and structural checks, not universal analytic correctness. Programs, input copies, exact commands, outputs, environment and hashes are retained here. Rerun commands are `python scripts/retain\_sources.py`, `python scripts/diagnostics.py`, `python scripts/build\_deliverables.py`, `python scripts/check\_deliverables.py`. TeX is saved without compilation; ROOT retains compilation/publication/acceptance and the protected editor. No supervised heavy-compute service was supplied, needed or used.

\end{document}
